\chapter[When Rules Constrain Us]{What If the Rules Themselves Constrain Us?}
\section{An Empty Cup}
Start by picking up an object: a cup.

Use the drinking cup beside you. Lift it and feel its weight. It has sides, a base, and perhaps a handle. Now for a question that sounds deliberately argumentative: which part of this cup is actually useful?

You will probably point to its sides: the ceramic, obviously; without that, the cup would be fragments on the floor. But think again---where does the water go? Into the ``emptiness'' enclosed by the sides. The sides are solid, but what lets you drink is precisely the space in the middle where there is nothing. A solid cup would be a lump of porcelain that could hold nothing.

Twenty-three centuries ago, someone turned this apparently contrary question into a key. He offered two further examples. Why can a wheel turn? Thirty spokes enter its hub, but the hub must be hollow for the axle to pass through. Why can people live in a house? Its walls enclose an empty space, while doors and windows are ``absences'' cut into them. His conclusion: solid things provide advantages, but empty things provide their actual use---``What is present offers benefit; what is absent provides use.''

This person was Laozi. At least, tradition attributes these words to him. The book he left contains only some five thousand Chinese characters. Called the \textit{Laozi}, it is more commonly known today as the \textit{Daodejing}. How much is five thousand characters? About three or four secondary-school essays. Yet for over two thousand years, people have copied, annotated, translated, and disputed this little book, which ranks among the world's most frequently translated works.

How could such a short book stir so much history? It asked a question nobody can escape, but few dare pursue to the end. We live in a net woven from rules---conventions, roles, laws, manners, marks, labels---originally meant to help us get through life. But \textbf{what if the rules themselves begin to constrain us?} This chapter follows that empty cup towards the answers offered by a group of Chinese thinkers over two millennia ago. Their answers are mischievous, full of fables, jokes, and argumentative twists, yet they pose one of the sharpest questions in the history of human thought.

\section{Cook Ding's Nineteen-Year-Old Knife}
Come into a scene from ``Nourishing Life'' in the \textit{Zhuangzi}, a Warring States text. Let us first watch it as an enlarged slow-motion sequence.

The setting is a kitchen: not your small kitchen at home, but a ruler's rear kitchen, with a chopping board the size of a bed. A freshly slaughtered ox lies on it, its hide gleaming, the smell of blood mingling with woodsmoke. Beside it stands its owner, Lord Wenhui, a ruler with his attendants. He has come to watch the rough work of butchery, perhaps merely while passing through.

The cook's surname is Ding, and people call him Cook Ding. He approaches with his sleeves rolled up. What follows makes everyone forget to blink.

A touch of the hand, a lean of the shoulder, a step of the foot, a press of the knee---the ox gives out a sequence of sounds as hide, flesh, and bone separate: swishing and rustling like rhythmic music. The knife travels through gaps between sinews and bones, almost too quickly to see, yet never once strikes bone. The watching ruler is astonished. How, he blurts out, can skill reach such heights?

Cook Ding lays down his knife and gives his famous reply. In essence: what I care about is the Dao, the Way, which goes beyond technique. When I first butchered oxen, I saw a whole ox---an enormous creature with nowhere to begin. After three years, I no longer saw a whole ox; I saw the structure of its sinews and bones and the paths between them. Now I do not look with my eyes at all. I meet it through the spirit: the senses stop and the spirit moves. The knife follows the body's natural grain, opening the gaps between sinews and bones and entering the hollows of the joints. These spaces belong to the ox's own structure, while my blade is so thin that it has almost no thickness. Put a blade without thickness into joints with space, and it has abundant room to move freely.

An ordinary cook, he says, replaces his knife every month because he hacks; a better cook replaces it every year because he cuts. But I have used this knife for nineteen years, butchered several thousand oxen, and its edge is still as fresh as if it had just left the whetstone.

Notice what is unusual here. Cook Ding does not disregard rules---quite the contrary, he knows the ``rules of the ox's body'' intimately. But he follows them differently from ordinary cooks. Others treat rules as obstacles and meet them head-on with their knives, ruining the blade and exhausting themselves. He understands them so thoroughly that he can move through them. For him, rules cease to be restraints and become paths.

When he has finished, Lord Wenhui says: Excellent! From hearing Cook Ding's words, I have learned how to nourish life.

An ox, a knife, a cook, a ruler. Keep this image in mind. Much of the rest of this chapter addresses its hidden question: facing the same structures, grain, and rules, why do some people wear their edges blunt while others keep theirs fresh for nineteen years?

\section{What Were Rules Originally For?}
Let us lay out the conflict before rushing to an answer.

First, we must acknowledge something: rules began as a well-intentioned invention. Imagine a school with no rules. Anyone can interrupt lessons, copy answers in exams, or barge into others in corridors. The weakest pupils would suffer first, because without rules, the strongest fists usually decide. Rules originally aimed to cage the strong and bring order to everyone. People in China over two millennia ago knew this well. Through centuries of the Spring and Autumn and Warring States periods, old conventions collapsed and rulers attacked one another: you attack me today, I attack you tomorrow, and ordinary farmers always suffer most.

When Confucius and his followers proposed rebuilding \textit{li}, ritual propriety---a whole system of roles, conduct, and duties---their starting point was entirely reasonable. Put everyone in their proper place: let fathers act as fathers, sons as sons, rulers as rulers, and ministers as ministers, and the machinery of society could work again. Rules were a net that settled people who otherwise collided in every direction.

Yet at the same time, another group looked at this net and saw something else. They asked: the net settles people, but does it not also entangle them?

Ritual propriety specifies how a son should stand, sit, and speak before his father; what ceremony a minister should perform before his ruler; how long funerals should last, which fabric mourners should wear, and how much they should weep. Regulation piles upon regulation, growing ever more detailed, until---how much of a person still belongs to them? From childhood someone hears ``You are a son,'' ``You are a minister,'' ``You are a student.'' Each identity comes with a complete script. If they perform it all their life, who has lived that life: the person or the script?

Worse, once established, rules grow by themselves. ``Respect your elders'' begins well, then grows into ``Elders are always right.'' ``Use examinations to select talent'' begins well, then becomes ``Marks define a person.'' Naming things begins as an aid to communication, but eventually names seem more real than the things themselves. Once someone is labelled a ``poor student,'' other people's view of them hardens, and so does their own.

A genuine problem emerges: \textbf{rules began as tools, so why do people become their servants? When rules constrain, distort, and swallow people, where can we step back to?}

I say ``step back'' because these thinkers often answer with retreat: back before rules grew, back before names were assigned, back into that useful empty cup. Later generations collectively called them ``Daoists.'' Their two central classics are the \textit{Daodejing} and the \textit{Zhuangzi}. But do not misunderstand: stepping back is not running away. We shall gradually see what it means.

Before continuing, choose a side. If a school rule is plainly unreasonable, would you say, ``Obey it first, whatever happens, then appeal through the proper procedure,'' or, ``An unreasonable rule loses its authority the moment it becomes unreasonable''? Remember your choice and see whether it has changed by the chapter's end.

\section{Two Strong Cases}
Now let us state both positions as fully as possible. We are not yet deciding who is right; first, we put both hands of cards on the table.

\textbf{Position One: Without ritual propriety, people fall apart. (The full case for rules.)}

Confucianism represents this position. Its reasoning deserves a fair hearing: for over two thousand years it has often been caricatured as ``stuffy old men making people kowtow,'' which is deeply unfair.

First: people are not born knowing how to be good. Everyone has desires. We want attractive things and seek revenge when hurt. Unless practised norms restrain and redirect these impulses from childhood, relationships become nothing but competition. Ritual propriety is like a river channel: without a channel, water is not freer; it floods disastrously.

Second: ceremonies and social roles are not merely appearances. You follow funeral rites not because ``the rules require crying,'' but because a solemn procedure helps a family find a place for immense grief. Calling a teacher ``teacher'' reminds both sides of their responsibilities. Names and ceremonies are invisible scaffolding, supporting what people cannot hold up inside themselves.

Third: abolishing rules is easy; clearing up afterwards is hard. Almost every historical celebration of ``smashing old conventions'' has been followed by more brutal power. The disappearance of rules does not mean equality has arrived; it means nobody can restrain the fist.

This is a strong case. Indeed, Chinese civilisation's continuity over thousands of years owes much to this technique of weaving norms into everyday life.

\textbf{Position Two: What is woven into everyday life may also be chains. (The Daoist challenge.)}

Laozi and Zhuangzi ask equally powerful questions. They do not deny the need for order; they question its cost and its origins.

First: do rules treat the disease or merely its symptoms? Why do people compete? Because they first distinguish ``This is good, that is bad,'' ``This is mine, that is yours.'' The finer the distinctions, the fiercer the struggle. Instead of removing its root, ritual propriety makes distinctions still more detailed: which vessels the high-ranking may use, which clothes the low-ranking must wear. Does this not map the battlefield more precisely?

Second: are performed conventions still genuine conventions? When bowing becomes compulsory, people learn to perform bows. Mourners can be hired for funerals; everyone at court can proclaim loyalty. The more complete the norms, the more accomplished the performance. Eventually society rewards not sincere people, but those who most resemble sincere people. The \textit{Daodejing} makes a cutting observation: only when the great Way is abandoned do people advocate benevolence and righteousness; only when families are in discord do they praise filial devotion and parental kindness. The louder a society shouts about virtue, the more it reveals virtue's scarcity.

Third: who made the rules? Daoists express this discreetly, yet it still stings two thousand years later. Rules never grow in a vacuum. Those who make them usually possess other kinds of power. When ``the rules'' become unquestionable, the greatest beneficiaries are often those best placed to interpret them.

Both cases stand up, and both are uncomfortable. Remember, too: Daoists are not ``troublemakers opposed to every rule,'' nor are Confucians ``bureaucratic accomplices in suppressing human nature.'' In actual history they argued for over two millennia and learned from each other for just as long. Many Chinese people carry both a Confucius and a Zhuangzi within them: Confucian when taking action, Daoist when wounded. This double aspect comes closer to reality than any label.

Widen the view and you will find a human-wide debate. In ancient Greece, Diogenes lived in a large wooden barrel. Alexander the Great supposedly visited and asked what favour he wanted; Diogenes merely asked him to move aside and stop blocking the sun. His whole life mocked the city-state's empty honours and courtesies. In nineteenth-century America, Thoreau built a cabin beside Walden Pond and lived there for over two years. In \textit{Walden}, he broadly records how schedules, possessions, and ``what others think'' drive people in circles; he wanted to strip life to its simplest form and see what remained. These people differ greatly, but share one gesture: stepping back to ask which conventions are genuinely necessary and which are merely familiar.

\section[Thinking Bridge: Change Your Position]{Thinking Bridge: Stand Somewhere Else}
Daoist books contain almost no dry argument; they overflow with stories, jokes, and contrariness. Not because their authors could not reason, but because their central tool is \textbf{changing perspective}: moving you from your position so that you can see how small it was. This tool divides into three moves, all available to you today.

\textbf{First move: exchange ``this'' for ``that.''}

``Equalising Things'' in the \textit{Zhuangzi} contains a famous idea. People constantly say ``This is right, that is wrong,'' yet ``this'' and ``that'' depend on each other. From this riverbank, the opposite bank is ``that''; cross over, and the original bank becomes ``that.'' Perspectives follow positions, but everyone treats their own position as the world's centre. Here is a simple exercise. Next time an argument makes you think ``Isn't it obvious?'', force yourself to describe the same issue from the other person's circumstances until they nod and say, ``Yes, that is what I mean.'' Until you can do that, you are opposing only an imaginary straw man.

\textbf{Second move: turn values upside down.}

``In the Human World'' in the \textit{Zhuangzi} describes a great tree whose twisted trunk and crooked branches cannot supply useful timber. Carpenters ignore it. Precisely because it is ``useless,'' it escapes the axe, grows enormously tall, and shelters passers-by in its shade. Zhuangzi asks through this story: you compete every day to be ``useful,'' yet useful things die of their usefulness. Fruit trees have branches broken during picking; good timber is felled for beams. What is use? What is uselessness? This reversal does not advise you to become worthless. It reminds you that ``useful'' and ``useless'' are verdicts under particular standards, and the standards deserve examination. A quiet child gifted at mathematics is useless under a standard that rewards outgoing personalities; under another standard, the same mind may be a treasure. Before discarding people or things judged ``useless,'' ask whose standard this is and what it measures.

\textbf{Third move: follow a paradox to its end.}

Daoist texts abound in paradoxes: apparently contradictory statements that open new perspectives on reflection. ``Great wisdom seems foolish,'' ``Great sound is scarcely heard,'' ``Act without forcing, yet leave nothing undone.'' What does a paradox do? Like a stick caught in gears, it brings habitual thinking to a sudden stop and forces another route. The ancient Greek Heraclitus liked this move too. The saying attributed to him---you cannot step into the same river twice---uses a knotty sentence to make you reconsider ``the same.'' Do not immediately dismiss a paradox as ``contradictory, therefore wrong.'' First ask: between which two familiar concepts does it want me to notice a previously unseen gap?

Together, these three moves are the mental exercises Daoism leaves us: change positions, reverse values, explore paradoxes. They do not hand you answers. Their one task is to loosen a screw or two in the frameworks welded into your mind. Only then can you tell which parts are load-bearing walls and which are merely partitions.

\section{Reading Passages from the Daodejing}
Now read the primary material. Take your time: every character in these passages has been disputed countless times.

The first passage opens the whole \textit{Daodejing}:

\begin{quote}
The Way that can be spoken is not the constant Way; the name that can be named is not the constant name.
\end{quote}
Literally: the ``Way'' expressible in language is not the enduring Way, and a ``name'' used for naming is not the enduring name. These twelve Chinese characters form the gateway to the book's five thousand. From the outset, it tells readers: the most important thing I am about to discuss is precisely what language cannot contain. You must pass through the words and reach for what lies behind them. Cling to the words, and you miss it.

The second passage comes from Chapter Eleven, the source of our opening cup:

\begin{quote}
Thirty spokes share one hub; in its absence lies the carriage's use. Knead clay to make a vessel; in its absence lies the vessel's use. Cut doors and windows to make a room; in its absence lies the room's use. Thus what is present offers benefit; what is absent provides use.
\end{quote}
In essence: thirty spokes converge on a hub, and its central hole lets the wheel work. Clay becomes a vessel, and its empty interior lets it hold things. Doors and windows are cut to make a house, and the space within its walls lets people inhabit it. Thus ``presence'' offers advantages, while ``absence'' is where usefulness lies.

Notice how to read this passage. On the surface are three everyday observations---a wheel, a clay cup, a house. Underneath is a carefully designed reasoning exercise: three examples you cannot doubt compel you to recognise a principle you have overlooked. Our entire civilisation grows on the side of ``having'': marks, certificates, property, reputation, packed schedules. Laozi points to the cup's emptiness and asks: have you counted that? An afternoon with nothing scheduled, a self defined by no role, a page read for no examination---these ``absences'' are precisely where things can be held.

Another passage worth knowing comes from Chapter Two:

\begin{quote}
When all under heaven recognise beauty as beauty, ugliness is already there; when all recognise goodness as goodness, what is not good is already there.
\end{quote}
Broadly: as soon as everyone knows what beauty is, the idea of ugliness arises with it; when everyone knows goodness, its opposite arises too. This does not mean ``Beauty is bad.'' It makes a conceptual point: beauty and ugliness, goodness and its opposite, difficulty and ease, height and lowness arise in pairs, like gears that define each other. You cannot take only one half. When a society frantically praises a particular kind of ``good,'' it simultaneously mass-produces people labelled ``bad.''

One final passage, from Chapter Twenty-Five, introduces another Daoist keyword: \textit{ziran}, ``naturalness.''

\begin{quote}
Humanity follows earth; earth follows heaven; heaven follows the Way; the Way follows what is so of itself.
\end{quote}
\textit{Fa} here means following or taking as a model. Humanity models itself on earth, earth on heaven, heaven on the Way, and the Way on ``naturalness.'' Be careful: this is not ``nature'' in today's sense of mountains, rivers, plants, and trees. In Classical Chinese, \textit{ziran} literally means ``so of itself'': not twisted by outside force or driven by prescription, everything grows and changes according to its own character. Vast as the Way is, its final model is those four Chinese characters, ``so of itself.'' This connects all the previous passages. Non-action lets things be so of themselves; the cup's ``absence'' makes room for this; rules become chains when they forbid people to be so of themselves, pruning every tree into one shape. Daoist ``following nature'' therefore does not mean lying in the mountains and neglecting everything. It means first understanding something's own nature, then treating it accordingly---an ox or a person alike.

Reading these three passages, you will notice that the \textit{Daodejing}'s manner of writing embodies its argument: brief, open, leaving room. It does not fill every space with words, because then you would remember only the words. It leaves emptiness for you to fill. We have returned to the cup.

\section{Three Readings of Non-Action}
We now reach the chapter's most delicate point. The \textit{Daodejing}'s central concept is \textit{wuwei}, ``non-action,'' perhaps the most badly misunderstood term in Chinese intellectual history. We now have enough pieces to set out three genuinely different readings and examine their reasons and limits.

\textbf{Reading One: Non-action means doing nothing.}

This is the most popular reading, and the one most in need of dismantling. On this account, non-action means opting out, giving up, and leaving things unattended. A student practising it does no homework; an official does no work; a parent leaves the children alone. There seems to be an argument: does ``non-action'' not literally mean ``not acting''?

But this reading breaks upon contact with the text. Chapter Forty-Eight states plainly:

\begin{quote}
In pursuing learning, increase daily; in pursuing the Way, decrease daily. Decrease and decrease again, until non-action is reached. Act without forcing, yet leave nothing undone.
\end{quote}
Notice the last phrase: non-action, yet \textbf{nothing left undone}. If non-action means doing nothing, how do you explain those three Chinese characters, \textit{wu bu wei}, ``nothing not done''? Recall Cook Ding: his blade moves freely and stays fresh for nineteen years. Can you say he has ``done nothing''? He has butchered thousands of oxen! His movements are quicker, more accurate, and more effective than anyone else's. Reading non-action as inactivity is not understanding Daoism; it is being deceived by literal wording. This is our ``easily misjudged counterexample'': Cook Ding looks relaxed and effortless, as though he has ``done nothing,'' yet he accomplishes more than anyone present.

\textbf{Reading Two: Non-action means avoiding reckless or forced intervention.}

This is a much more serious reading. In Classical Chinese, \textit{wei} means not only ``doing'' but also carries the sense of deliberate, forcing action. Non-action means not working forcibly against the grain of things. Anyone who grows flowers understands: you cannot help seedlings grow by pulling them upwards. The farmer in the \textit{Mencius} who does just that exemplifies ``too much action.'' Managing a class is similar. If the teacher controls every detail, students can only comply or resist. Leave space, and order may grow by itself.

This reading has strong support. It explains ``non-action, yet nothing left undone'': refrain from meddling, and things can succeed. It also matches the ``less is more'' experience of modern education and management. But its limitation is clear: where is the boundary? What separates ``non-intervention'' from ``irresponsibility''? When a child is absorbed in gaming, is parental non-action respect for how things work or neglect of duty? This reading supplies a principle, not a measuring scale.

\textbf{Reading Three: Non-action is a state of highly developed skill.}

This reading puts Cook Ding at the centre. Non-action is not a beginner's inactivity but the ``no need to force it'' that comes with consummate mastery: practising an instrument until the fingers move by themselves, a ball game until the body arrives before conscious thought, a knife until it follows an ox's gaps. The person feels no ``doing,'' yet does it more beautifully than anyone. Person and rule cease to be separate. Rules are completely internalised as intuition, and intuition becomes rule. In today's terms, this approaches psychological ``flow,'' the loss of self-awareness in total absorption.

This reading explains the most: Cook Ding, ``nothing left undone,'' and Daoism's repeated emphasis on ``forgetting''---technique, conventions, the self. Its limitation is that non-action becomes a peak experience for a few. What, then, does Daoism offer ordinary people and daily life? And how can a butcher's skilled state extend to governing a country or arranging a life? Between the kitchen and the world, the \textit{Zhuangzi} supplies no map.

All three readings hold part of the truth. Reading One is a trap, so we exclude it; Two and Three need not be alternatives. One says ``Do not meddle''; the other, ``Practise until mastery transforms you.'' But notice what we have done: faced with one term, we split ``What does it mean?'' into competing explanations, each with blind spots. This move matters more than choosing an answer. Next time someone online declares ``Daoism just tells people to give up,'' you have tools to take that claim apart.

\section[Further Challenge: Beyond Language]{Further Challenge: What Language Cannot Contain}
Now for the chapter's weightiest theoretical problem, hidden in the \textit{Daodejing}'s opening twelve characters: ``The Way that can be spoken is not the constant Way; the name that can be named is not the constant name.'' Why can language not contain what matters most? Is this mystification or a serious philosophical discovery?

Consider the famous story in ``The Way of Heaven'' in the \textit{Zhuangzi}. Duke Huan of Qi is reading in his hall. Below him works an old wheelwright called Bian. Bian puts down his chisel, approaches, and asks what the ruler is reading. The words of sages, says the duke. Are the sages alive? No, long dead. Bian delivers a startling verdict: then what you are reading is merely the dregs of the ancients.

The duke is furious. A wheelmaker dares judge his reading? Explain yourself, or die.

Unruffled, Bian uses his craft as an analogy. In essence: when shaping a wheel, make the joint too loose and it will not hold; too tight and it will not fit. Getting it neither loose nor tight means finding it in the hand and responding in the heart. The hand feels it, the heart responds, but the mouth cannot express it. I cannot even teach this judgement to my son, and he cannot receive it from me, so at seventy I still have to do the work myself. The ancient sages and what they truly could not express have died together. What you read, then, can only be their dregs.

This story is frighteningly sharp, but not anti-intellectual. Bian does not say ``All books are rubbish.'' He says: \textbf{some knowledge cannot be compressed into language for transport.} Modern philosophers call this ``tacit knowledge'': you can know more than you can say. Cycling is a standard example. You may be an excellent cyclist yet unable to write a single instruction that teaches someone else to ride. How do you balance? Which way do you lean, and how far? You ``know'' all this, but the knowledge lives in your muscles, not a dictionary. Swimming, wine tasting, noticing slight changes in a friend's mood, and a ``feel'' for mathematical problems all belong here.

Now ``the name that can be named is not the constant name'' looks less mysterious. What is language? A classification system: it cuts a continuous world into compartments and labels each---``table,'' ``chair,'' ``right,'' ``wrong,'' ``good student,'' ``poor student.'' Classification is useful; without it we could hardly move. But Daoists ask us to inspect the bill. First, humans choose where to cut; the divisions do not come built into the world. The same cloud can be ``a cumulonimbus'' or ``the one that looks like a rabbit.'' Second, once compartments are labelled, people mistake labels for actual things and forget they made the divisions themselves. Third, and most importantly, the experiences that matter most are hardest to fit into compartments: inexpressible sadness, inexplicable affection, Cook Ding's judgement in his hands, Bian's sense of tightness and looseness.

This problem is not uniquely Chinese. The twentieth-century Austrian philosopher Wittgenstein used an exceptionally rigorous little book to trace the logic of language. At its conclusion, he said in effect: where speech is impossible, we must remain silent. Interestingly, this can express disdain---do not discuss what you cannot state clearly---or reverence: precisely because what matters most cannot be said, we must be humble before language. Daoists would choose the latter. The famous exchange above the Hao River in ``Autumn Floods'' in the \textit{Zhuangzi} plays with the same problem. Watching fish, Zhuangzi says: how freely they swim; that is their happiness. His friend Huizi challenges him: you are not a fish, so how do you know its happiness? Zhuangzi replies: you are not me, so how do you know I do not know its happiness? Often treated as verbal sparring, this exchange probes a deep question: \textbf{what grounds all our ``knowledge'' of other people's inner lives?} Language transports propositions, not experiences. But if we conclude that ``people cannot understand one another at all,'' what are we doing in this exchange right now?

This section offers no answer, only a map. Language is an excellent tool and a transparent wall. Seeing that wall---knowing that each word leaves something out---does not mean falling silent. It means leaving room, in speaking and listening, for what cannot be said. A cup holds water because it is empty; thought holds the world by recognising that language cannot fill it completely.

\section{Check-Ins, Labels, and Algorithms}
This question, over two thousand years old, now sits inside your phone.

Start with a familiar device: the daily check-in. Vocabulary check-ins, running check-ins, reading check-ins. Originally, it is a gentle tool: divide a large goal into small boxes and let the achievement of ``N consecutive days'' encourage you. After a while, however, many people notice a subtle drift. Somehow ``Did I read today?'' gives way to ``Did I check in today?'' Two pages may suffice, but the streak must continue; shaking a phone may manufacture steps, but missing a day is unbearable. The goal quietly changes from ``becoming a better person'' to ``maintaining the record itself.'' In Daoist terms, names once served reality; now reality serves names. What would Laozi say about someone remembering at 11:30 p.m. that they have not checked in, climbing out of bed, and turning two pages? He might point to the empty cup: a cup exists to hold water, not to prove that it is a cup.

Now labels. Your age enjoys a bumper harvest of them. Personalities have four-letter codes, interests long strings of tags, and music apps announce ``You are this kind of person'' more confidently than you do. Classification is not wrong; without it the world is unmanageable. But remember Chapter Two's paired gears: every category of ``this kind of person'' creates ``not this kind of person.'' When a sixteen-year-old earnestly says, ``I'm an I-type, an introvert, so I cannot speak in class,'' gently ask: does the label help you understand yourself, or decide who you are for you? A label should be a shoe to try on, not a shackle welded to your ankle. This is why Daoists distrust fixed classifications: names can cultivate the reality they name.

Then there is the largest classification machine: the algorithm. Recommendation systems observe what you watch, how long you linger, and what you click, drawing an increasingly detailed portrait before pruning the world into what that portrait likes. Notice the structure: it uses pure ``presence.'' Every click and second of attention is data. What you did not click, your inexplicable melancholy, the ten minutes you spent gazing out of the window this afternoon---these count as zero in its ledger, as nonexistent. Yet Zhuangzi might say that something essential was happening during those ten idle minutes: the empty part of the cup. The algorithm is not malicious; it simply cannot see ``absence.'' Someone raised by a system that rewards only ``presence'' must find ways to defend absences: unrecorded reading, aimless walks, feelings never posted.

A further ancient suspicion sounds prophetic today. Chapter Twelve says: ``The five colours blind the eyes; the five tones deafen the ears; the five flavours dull the palate.'' Broadly, excessive bombardment by colours, sounds, and tastes disables the senses. Here lies Daoist suspicion of desire. The concern is not that people have desires, but that desires can be \textbf{fed}: stronger stimulation demands a stronger next dose. Like a hamster on a wheel, the harder a person runs, the less they know where they are going. Look at your phone: videos become ever more explosive, game levels more exciting, shopping apps know what you want before you do. Their engineers may never have read the \textit{Daodejing}, but their work illustrates the warning in reverse: studying how to make colours more dazzling and sounds more piercing, keeping your ``wanting'' forever just short of satisfaction. The Daoist response is not abstinence, but clear accounting: did this desire grow inside me, or did someone plant it there?

Consider, too, the fashionable Chinese terms \textit{neijuan}, exhausting inward competition, and \textit{tangping}, ``lying flat.'' Competition becomes white-hot, and some announce their withdrawal: no more striving; lie flat. You should now recognise these as opposite ends of one seesaw: being swallowed by rules or overturning the entire rule system. Cook Ding offers a third path. He neither strikes bone head-on nor throws down his knife and leaves; he learns the grain well enough to move through it. Non-action is not lying flat. Lying flat is Reading One's modern version, already dismantled against the original text. The real skill is seeing this game's structure---which sinews and bones must be negotiated, which imposing names merely intimidate---then taking the thinnest blade through the widest space. Today this has a plain name: ``Do not mistake someone else's racecourse for your life.''

Finally, a quieter echo. Daoist suspicion of technology appears in Chapter Eighty's picture of small states with few people. They possess various devices but do not use them. Neighbours are close enough to hear one another's chickens and dogs, yet people remain content with their own days, growing old without visiting and disturbing one another. For two millennia, many mocked this as turning the clock back. Today, when hundreds of millions voluntarily practise ``digital fasting,'' delete apps, or buy old phones that only make calls, the laughter sounds less assured. The question was never simply whether to have technology, but who serves whom. Do objects revolve around the empty space within you, or do you revolve around objects until nothing remains inside?

\section[Your Turn: Axe or Cook Ding's Knife?]{Your Turn: Chop Down the Tree, or Learn from Cook Ding?}
Return to the opening question and push it to its extreme.

Suppose you chair the student council and hold a list of long-standing grievances: rigid schedules, pointless ceremonies, assessment based solely on marks, endless check-in statistics. Students divide into two camps. One says: rules are chains; abolish them all and let everyone grow freely. The other says: rules are order; loosen none, because one breach brings down the whole embankment.

Both camps quote this chapter. The abolitionists recite ``The sage is not benevolent'' and tell of the tree that survived by being useless. The defenders accuse Daoists of misreading rules and recall the chaos that followed the destruction of old conventions.

But you now possess more than slogans. You know rules began with good intentions, and that they grow by themselves and can swallow people. You know non-action is not inactivity: Cook Ding's knife remained fresh through something other than going on strike. You know language deceives, and that large words like ``freedom'' and ``order'' conceal things left unsaid. You also know the tool of changing positions: first cross over and state the other side's full case, then return to make your cut.

So answer this: facing a constraining net of rules, what is the right move? Swing an axe like a woodcutter and chop through it, or learn its structure like Cook Ding and find the gaps? The net-cutter may become a new net-maker who binds others---history offers plenty of examples. The person slipping through may grow accustomed to the gaps and forget that the net need not exist at all---there are still more such people.

In that particular school, which rules should go, which should stay, and which should remain but be enforced differently? What is your criterion: ``Is it useful?'' or ``Whom does it turn into a servant?''

There is no standard answer. But give an answer and reasons. As you write those reasons, watch every large word you use---``freedom,'' ``growth,'' ``rules,'' ``prospects.'' Each is a compartment cut by someone else. You may use them, but from today you should know that you stand inside compartments, and that beyond them there remains a whole cupful of emptiness.
